\textbf{Phần 1. Trắc nghiệm lựa chọn}
\setcounterpageref{ex}{0}

\begin{ex} %[2D7V36]
Tính giá trị của biểu thức: $72 - 38 + 26$.
\choice
{$58$}
{\True $60$}
{$59$}
{$61$}
\loigiai{
$72 - 38 = 34$; $34 + 26 = 60$.
\textbf{Đáp án của câu hỏi là B}
}
\end{ex}

\begin{ex} %[2D7V36]
Một lớp học có $36$ học sinh. Trong đó có $19$ bạn nữ. Số bạn nam của lớp là
\choice
{\True $17$ bạn}
{$16$ bạn}
{$18$ bạn}
{$15$ bạn}
\loigiai{
Số bạn nam: $36 - 19 = 17$ (bạn).
\textbf{Đáp án của câu hỏi là A}
}
\end{ex}

\begin{ex} %[2D7V36]
Số bé nhất có hai chữ số mà tổng hai chữ số của nó bằng $9$ là
\choice
{$27$}
{$36$}
{\True $18$}
{$90$}
\loigiai{
Chữ số hàng chục nhỏ nhất là 1, hàng đơn vị là $9 - 1 = 8$. Số đó là 18.
\textbf{Đáp án của câu hỏi là C}
}
\end{ex}

\begin{ex} %[2D7V36]
Chu vi hình tứ giác có độ dài các cạnh lần lượt là $14\text{ cm}, 18\text{ cm}, 22\text{ cm}, 26\text{ cm}$ là
\choice
{$78\text{ cm}$}
{$79\text{ cm}$}
{$82\text{ cm}$}
{\True $80\text{ cm}$}
\loigiai{
$(14 + 26) + (18 + 22) = 40 + 40 = 80\text{ (cm)}$.
\textbf{Đáp án của câu hỏi là D}
}
\end{ex}

\begin{ex} %[2D7V36]
Can thứ nhất có $18\text{ L}$ dầu, can thứ hai có nhiều hơn can thứ nhất $6\text{ L}$ dầu. Cả hai can có bao nhiêu lít dầu?
\choice
{$40\text{ L}$}
{\True $42\text{ L}$}
{$44\text{ L}$}
{$24\text{ L}$}
\loigiai{
Can thứ hai có: $18 + 6 = 24\text{ (L)}$. Cả hai can có: $18 + 24 = 42\text{ (L)}$.
\textbf{Đáp án của câu hỏi là B}
}
\end{ex}

\begin{ex} %[2D7V36]
Tìm $x$ biết: $x + 27 = 83 - 19$.
\choice
{\True $37$}
{$35$}
{$36$}
{$38$}
\loigiai{
Vế phải: $83 - 19 = 64$. Ta có: $x = 64 - 27 = 37$.
\textbf{Đáp án của câu hỏi là A}
}
\end{ex}

\begin{ex} %[2D7V36]
Nếu thứ Bảy tuần này là ngày $16$ thì thứ Bảy tuần sau là ngày mấy?
\choice
{$22$}
{$24$}
{\True $23$}
{$25$}
\loigiai{
$16 + 7 = 23$.
\textbf{Đáp án của câu hỏi là C}
}
\end{ex}

\begin{ex} %[2D7V36]
Một sợi dây dài $65\text{ dm}$. Cắt bớt đi $28\text{ dm}$. Chiều dài đoạn dây còn lại là
\choice
{$35\text{ dm}$}
{$36\text{ dm}$}
{$38\text{ dm}$}
{\True $37\text{ dm}$}
\loigiai{
$65 - 28 = 37\text{ (dm)}$.
\textbf{Đáp án của câu hỏi là D}
}
\end{ex}

\begin{ex} %[2D7V36]
Tìm một số biết rằng lấy số đó trừ đi $35$ rồi cộng với $18$ thì được $60$.
\choice
{$75$}
{\True $77$}
{$76$}
{$78$}
\loigiai{
Tính ngược từ cuối: $60 - 18 + 35 = 42 + 35 = 77$.
\textbf{Đáp án của câu hỏi là B}
}
\end{ex}

\begin{ex} %[2D7V36]
Trong một phép trừ, số bị trừ là $90$, số trừ là $47$. Hiệu của phép trừ đó là
\choice
{\True $43$}
{$41$}
{$42$}
{$44$}
\loigiai{
$90 - 47 = 43$.
\textbf{Đáp án của câu hỏi là A}
}
\end{ex}

\begin{ex} %[2D7V36]
Đồng hồ chỉ $20$ giờ $15$ phút, tức là mấy giờ đêm?
\choice
{$7$ giờ $15$ phút tối}
{$9$ giờ $15$ phút tối}
{\True $8$ giờ $15$ phút tối}
{$8$ giờ rưỡi tối}
\loigiai{
$20$ giờ $= 8$ giờ tối. Vậy là $8$ giờ $15$ phút tối.
\textbf{Đáp án của câu hỏi là C}
}
\end{ex}

\begin{ex} %[2D7V36]
Bác Nam có $62\text{ kg}$ thóc, bác bán đi $38\text{ kg}$ thóc. Bác Nam còn lại bao nhiêu ki-lô-gam thóc?
\choice
{$22\text{ kg}$}
{$23\text{ kg}$}
{$25\text{ kg}$}
{\True $24\text{ kg}$}
\loigiai{
$62 - 38 = 24\text{ (kg)}$.
\textbf{Đáp án của câu hỏi là D}
}
\end{ex}

\textbf{Phần 2. Trắc nghiệm đúng sai}
\setcounterpageref{ex}{0}

\begin{ex} %[2D7V36]
Xét tính đúng sai của các bài toán tổng hợp toàn bộ Học kì 1:
\choiceTF
{\True Tổng của số lẻ lớn nhất có 2 chữ số và số chẵn nhỏ nhất có 2 chữ số là 109.}
{Đoạn thẳng dài 1 dm dài hơn đoạn thẳng dài 12 cm.}
{\True Trong một phép trừ, nếu số bị trừ bằng số trừ thì hiệu bằng 0.}
{Một năm chỉ có 3 tháng có 31 ngày.}
\loigiai{
\begin{itemize}
    \item Ý a) Đúng: $99 + 10 = 109$. Đúng.
    \item Ý b) Sai: $1\text{ dm} = 10\text{ cm} < 12\text{ cm}$. Sai.
    \item Ý c) Đúng: Đúng.
    \item Ý d) Sai: Có tất cả 7 tháng có 31 ngày.
\end{itemize}
\begin{itemize}
    \item a) \textbf{Đúng}.
    \item b) \textbf{Sai}.
    \item c) \textbf{Đúng}.
    \item d) \textbf{Sai}.
\end{itemize}
}
\end{ex}

\begin{ex} %[2D7V36]
Một giá sách có hai tầng, tầng trên có 48 quyển, tầng dưới có 36 quyển:
\choiceTF
{\True Cả hai tầng có tất cả 84 quyển sách.}
{\True Tầng trên nhiều hơn tầng dưới 12 quyển sách.}
{Chuyển 12 quyển từ tầng trên xuống tầng dưới thì hai tầng bằng nhau.}
{\True Nếu chuyển 6 quyển từ tầng trên xuống tầng dưới thì mỗi tầng có 42 quyển.}
\loigiai{
\begin{itemize}
    \item Ý a) Đúng: $48 + 36 = 84$ quyển.
    \item Ý b) Đúng: $48 - 36 = 12$ quyển.
    \item Ý c) Sai: Chỉ cần chuyển $12 : 2 = 6$ quyển là bằng nhau.
    \item Ý d) Đúng: $48 - 6 = 42$ và $36 + 6 = 42$. Đúng.
\end{itemize}
\begin{itemize}
    \item a) \textbf{Đúng}.
    \item b) \textbf{Đúng}.
    \item c) \textbf{Sai}.
    \item d) \textbf{Đúng}.
\end{itemize}
}
\end{ex}

\textbf{Phần 3. Trả lời ngắn}
\setcounterpageref{ex}{0}

\begin{ex} %[2D7V36]
Tính giá trị biểu thức: $48 + 36 - 29$.
\shortans{55}
\loigiai{
\textbf{Đáp số:} $55$\\
$48 + 36 = 84$; $84 - 29 = 55$.
}
\end{ex}

\begin{ex} %[2D7V36]
Một đàn gà có 63 con, người ta bán đi 37 con. Hỏi đàn gà còn lại bao nhiêu con?
\shortans{26}
\loigiai{
\textbf{Đáp số:} $26$\\
$63 - 37 = 26$ con.
}
\end{ex}

\begin{ex} %[2D7V36]
Tìm x biết: $85 - x = 27 + 19$.
\shortans{39}
\loigiai{
\textbf{Đáp số:} $39$\\
$27 + 19 = 46 \Rightarrow x = 85 - 46 = 39$.
}
\end{ex}

\begin{ex} %[2D7V36]
Tính chu vi hình tứ giác có độ dài các cạnh là 15 cm, 25 cm, 18 cm, 22 cm.
\shortans{80}
\loigiai{
\textbf{Đáp số:} $80$\\
$(15 + 25) + (18 + 22) = 40 + 40 = 80\text{ cm}$.
}
\end{ex}

\begin{ex} %[2D7V36]
Can A có 35 L dầu, Can B có ít hơn Can A 8 L dầu. Hỏi cả hai can có bao nhiêu lít dầu?
\shortans{62}
\loigiai{
\textbf{Đáp số:} $62$\\
Can B: $35 - 8 = 27\text{ L}$. Cả hai can: $35 + 27 = 62\text{ L}$.
}
\end{ex}

\begin{ex} %[2D7V36]
Nếu hôm nay là thứ Tư ngày 7 thì thứ Tư của 2 tuần sau là ngày mấy?
\shortans{21}
\loigiai{
\textbf{Đáp số:} $21$\\
$7 + 7 + 7 = 21$.
}
\end{ex}
